\begin{center}
\begin{tikzpicture}[x=.88cm,y=.88cm,font=\small]
  \fill[paper] (0,0) rectangle (10.5,5.7);
  \draw[blue!55,line width=1pt] (.2,4.8) .. controls (3.2,4.6) and (6.3,5.1) .. (10.0,4.6);
  \node[text=blue!70!black,above] at (5.0,4.9) {黄河};
  \draw[blue!45,line width=.8pt] (3.8,.4) .. controls (5.0,2.0) and (6.2,2.2) .. (8.8,.7);
  \node[text=blue!70!black,below] at (6.2,1.35) {淮水方向};

  \fill[green!20] (1.0,3.5) -- (4.0,3.3) -- (4.2,5.0) -- (1.6,5.2) -- cycle;
  \node[align=center] at (2.5,4.25) {晋\\北方联盟};
  \fill[purple!18] (4.8,.8) -- (8.4,.9) -- (8.7,3.4) -- (5.5,3.6) -- cycle;
  \node[align=center] at (6.7,2.2) {楚\\北上军队};
  \fill[yellow!25] (7.6,3.5) -- (9.4,3.4) -- (9.6,4.7) -- (7.9,4.9) -- cycle;
  \node at (8.6,4.1) {齐};
  \fill[orange!24] (4.1,3.2) -- (5.1,3.2) -- (5.2,4.1) -- (4.2,4.2) -- cycle;
  \node at (4.65,3.65) {宋};
  \fill[timeline] (5.1,3.0) circle (.14);
  \node[below,align=center] at (5.1,2.9) {城濮\\战场};

  \draw[-{Stealth[length=2mm]},ultra thick,color=green!55!black] (2.8,3.75) -- (4.95,3.05);
  \draw[-{Stealth[length=2mm]},ultra thick,color=dynasty] (6.1,2.6) -- (5.25,2.95);
  \node[text=green!55!black,above] at (3.75,3.35) {晋及盟军};
  \node[text=dynasty,below] at (5.65,2.55) {楚军北上};
\end{tikzpicture}

\smallskip
\small\textit{图  城濮战役空间示意：突出晋的北方联盟、楚的北上方向、宋等中原国家和战场的相对位置；不表示精确行军路线或疆界。}
\end{center}
